\subsection{KG Storage Model}

The Knowledge Graph stores triples using dictionary encoding and SQLite persistence.

\subsubsection{Dictionary Encoding}

String values are mapped to compact integer IDs to reduce storage overhead and enable fast comparison:

\begin{itemize}
    \item \textbf{Entity dictionary}: $(entity\_id : \mathbb{N}) \mapsto uri : \texttt{str}$
    \item \textbf{Predicate dictionary}: $(pred\_id : \mathbb{N}) \mapsto uri : \texttt{str}$
    \item \textbf{Literal table}: $(lit\_id : \mathbb{N}) \mapsto (value : \texttt{str}, datatype : \texttt{str})$
\end{itemize}

\subsubsection{Triple Storage}

The core storage schema is:

\begin{verbatim}
CREATE TABLE triples (
    subject    INTEGER NOT NULL,
    predicate  INTEGER NOT NULL,
    object     INTEGER NOT NULL,
    obj_type   TEXT NOT NULL,  -- 'entity' or 'literal'
    ev_id      INTEGER DEFAULT 0
);
\end{verbatim}

Three primary indexes are maintained:
\begin{verbatim}
CREATE INDEX idx_spo ON triples(subject, predicate, object);
CREATE INDEX idx_pos ON triples(predicate, object, subject);
CREATE INDEX idx_ops ON triples(object, predicate, subject);
\end{verbatim}

\subsubsection{Reification Schema}

For $n$-ary facts, reification connects roles and values through an event node:

\begin{verbatim}
CREATE TABLE reification (
    ev_id   INTEGER NOT NULL,  -- event node ID
    pred_id INTEGER NOT NULL,  -- role predicate ID
    obj_id  INTEGER NOT NULL,  -- role value ID (entity or literal)
    role    TEXT NOT NULL       -- role name
);
\end{verbatim}

\subsubsection{Persistence}

All data is persisted to SQLite with WAL journaling for concurrent read access. On load, all indexes are automatically rebuilt. Export produces N-Triples format; import parses N-Triples back into the internal representation.
